% bank_fintech_substance/drafts/zh/manuscript.tex
% P3 中文 manuscript 骨架 — M0 G0 scaffolded per B-route plan v2.55 (2026-05-17)
% 编译: xelatex + biber + xelatex × 2

% papers_assets/_shared/preamble_zh.tex
% Shared Chinese LaTeX preamble (XeLaTeX + CJKutf8).
% P4.1 (plan-r3 derivative-paper extension, 2026-05-15)
% Compile: xelatex + biber + xelatex × 2

\documentclass[12pt,a4paper]{article}

% --- CJK Fonts (XeLaTeX, cross-platform Q1.1b fallback) ---
\usepackage{fontspec}
\usepackage{xeCJK}
% Windows (local dev): SimSun/SimHei/FangSong
% Linux/macOS (CI): Noto Serif CJK SC / Noto Sans CJK SC (fonts-noto-cjk apt package)
% Last resort: AR PL UMing CN (fonts-arphic-uming) which is preinstalled in many texlive distros
\IfFontExistsTF{SimSun}{%
  \setCJKmainfont{SimSun}[BoldFont=SimHei]%
}{%
  \IfFontExistsTF{Noto Serif CJK SC}{%
    \setCJKmainfont{Noto Serif CJK SC}[BoldFont=Noto Serif CJK SC]%
  }{%
    \setCJKmainfont{AR PL UMing CN}[AutoFakeBold=2]%
  }%
}
\IfFontExistsTF{SimHei}{%
  \setCJKsansfont{SimHei}%
}{%
  \IfFontExistsTF{Noto Sans CJK SC}{%
    \setCJKsansfont{Noto Sans CJK SC}%
  }{%
    \setCJKsansfont{AR PL UKai CN}%
  }%
}
\IfFontExistsTF{FangSong}{%
  \setCJKmonofont{FangSong}%
}{%
  \IfFontExistsTF{Noto Sans Mono CJK SC}{%
    \setCJKmonofont{Noto Sans Mono CJK SC}%
  }{%
    \setCJKmonofont{AR PL UMing CN}%
  }%
}
\setmainfont{Times New Roman}[Ligatures=TeX]
\xeCJKsetup{CJKecglue={\hskip 0.15em plus 0.05em minus 0.05em}}

% --- Math ---
\usepackage{amsmath, amssymb, amsthm}

% --- Tables & Figures ---
\usepackage{booktabs}
\usepackage{graphicx}
\usepackage{longtable}
\usepackage{rotating}
\usepackage{caption}
\captionsetup{font=small,labelfont=bf,skip=4pt}

% --- Layout ---
\usepackage[margin=1in]{geometry}
\usepackage{setspace}
\onehalfspacing
\usepackage{hyperref}
\hypersetup{colorlinks=true,linkcolor=blue,citecolor=blue,urlcolor=blue}

% --- Bibliography ---
\usepackage[backend=biber,style=apa,natbib=true]{biblatex}

% --- Macros ---
\newcommand{\code}[1]{\texttt{#1}}
\usepackage[dvipsnames]{xcolor}
\newcommand{\NoteThis}[1]{\textbf{\textcolor{red}{[TODO: #1]}}}

% --- Theorem environments ---
\newtheorem{prop}{命题}
\newtheorem{hyp}{假设}


\addbibresource{refs.bib}

\title{实质性 vs 象征性金融科技战略：\\来自中国银行业长期 ROA 反转的证据}
\author{\NoteThis{作者姓名}}
\date{\today}

\begin{document}
\maketitle

\begin{abstract}
\NoteThis{最终摘要待定 — 目标约 300 中文字。}
本文将 Westphal--Zajac (1994, 1998 \emph{ASQ}) 的实质性 vs 象征性管理框架应用于
2007--2024 年间的银行 fintech 战略宣告。通过大语言模型对每个银行—年度的 fintech
公告进行二元 substantiveness 分类 (\code{innovation\_substantiveness\_full}), 检验
W\&Z 隐含的 dual outcome 预测: 象征性宣告应产生正向短期市场反应但长期 ROA 反转
(within-bank 3 年前向变化), 实质性宣告则不应。识别策略利用 4 个外生政策冲击
(余额宝上线 / 网络借贷整顿 / 资管新规 / 蚂蚁 IPO 暂停) 做截面强度 DiD。研究发现
\NoteThis{核心结论}\@。
\end{abstract}

\section{引言}\label{sec:intro}
\NoteThis{研究问题动机与贡献。}

\section{文献综述}\label{sec:lit}
\NoteThis{三条文献: (1) Westphal--Zajac 实质性/象征性
\autocite{westphal1994substance,westphal1998sym,westphal2001decoupling};
(2) 战略变革的象征管理 \autocite{fisszajac2006symbolic};
(3) Practice-based view \autocite{bromiley2014practice}。}

\section{方法}\label{sec:method}
\subsection{样本与数据}
银行—年度面板覆盖 $N \times T \approx 3000\text{--}9000$ 银行—年, 2007--2024
(per spec.sample), 经 \code{news\_process} 流水线 (stage 32a innovation\_substantiveness)。

\subsection{实质性度量}
\NoteThis{LLM 标注流水线 (innovation\_substantiveness\_full 二元指示); full vs ctx
双轨交叉验证。}

\subsection{Dual Outcome 识别策略}
\NoteThis{Long\_run\_roa\_change\_3yr (within-bank 3 年前向) + car\_3day\_window
\_fintech\_announcement (短期 CAR); within-bank within-year FE; bank-level cluster。}

\section{数据}\label{sec:data}
\NoteThis{数据来源、样本选择、描述统计 (Table 1)。}

\section{实证结果}\label{sec:results}
\NoteThis{Table 2 (M1$\to$M6 渐进主回归); Table 3 (dual outcome 短期 CAR + 长期 ROA
反转联合); Table 4 (异质性 4 银行类型 $\times$ SOE $\times$ wave $\times$ 资本充足率
quartile); Table 5 (稳健性 alternative substantiveness track full vs ctx)。}

\section{结论}\label{sec:conclusion}
\NoteThis{总结发现与战略管理理论启示。}

\printbibliography

\end{document}
